% Part: propositional-logic
% Chapter: syntax-and-semantics
% Section: introduction

\documentclass[../../../include/open-logic-section]{subfiles}

\begin{document}

\olfileid[pt]{pl}{syn}{int}

\olsection{Introdução}

A lógica proposicional trata de !!{formula}s construídas a partir de
!!{propositional variable}s por meio dos conectivos proposicionais
$\lnot$, $\land$, $\lor$, $\lif$ e $\liff$. Intuitivamente,
!!a{propositional variable}~$p$ representa uma sentença ou proposição
que é verdadeira ou falsa. Sempre que se determina o ``valor de verdade'' da
!!{propositional variable} em !!a{formula}, determina-se também o
valor de verdade de quaisquer !!{formula}s formadas a partir dessas variáveis por meio de conectivos
proposicionais. Dizemos que a lógica proposicional é \emph{verofuncional},
pois sua semântica é dada por funções de valores de
verdade. Em particular, na lógica proposicional deixamos de
considerar qualquer determinação adicional da verdade e da falsidade, por exemplo,
se algo é necessariamente verdadeiro, em vez de apenas contingentemente
verdadeiro; se algo é conhecido como verdadeiro; ou se algo
é verdadeiro agora, em vez de ter sido ou vir a ser verdadeiro. Consideramos apenas
dois valores de verdade: verdadeiro ($\True$) e falso ($\False$); assim, excluímos
da discussão a possibilidade de uma afirmação não ser nem verdadeira
nem falsa, ou ser apenas parcialmente verdadeira. Também nos concentramos apenas nos conectivos para os quais
o valor de verdade de !!a{formula} construída a partir deles é completamente
determinado pelos valores de verdade de suas partes (e não, por exemplo, por seu
significado). Em particular, é controverso se o valor de verdade dos condicionais em
inglês é verofuncional nesse sentido. O condicional material~$\lif$ é;
outras lógicas tratam de condicionais que não são verofuncionais.

Para desenvolver a teoria e a metateoria da lógica proposicional
verofuncional, devemos primeiro definir a sintaxe e a semântica de
suas expressões. Descreveremos uma maneira de construir
!!{formula}s a partir de !!{propositional variable}s usando os conectivos.
Definições alternativas são possíveis. Outros sistemas escolhem
símbolos diferentes, selecionam conjuntos diferentes de conectivos como
primitivos e usam parênteses de maneira diferente (ou nem os usam,
como no caso da chamada notação polonesa). O que todas as abordagens
têm em comum, porém, é que as regras de formação definem o conjunto de
!!{formula}s \emph{indutivamente}. Se isso for feito corretamente, toda expressão
poderá resultar, em essência, de uma única maneira segundo as regras de
formação. A definição indutiva que produz expressões de
\emph{leitura unívoca} significa que podemos atribuir significado a essas
expressões pelo mesmo método---uma definição indutiva.

Atribuir significado às expressões é o domínio da semântica. O
conceito central da semântica da lógica proposicional é o de
satisfação em !!a{valuation}. !!^a{valuation}~$\pAssign{v}$ atribui
os valores de verdade $\True$, $\False$ às !!{propositional variable}s. Toda
!!{valuation} determina um valor de verdade $\pValue{v}(!A)$ para qualquer
!!{formula}~$!A$. !!^a{formula} é satisfeita em
!!a{valuation}~$\pAssign{v}$ se, e somente se, $\pValue{v}(!A) = \True$---escrevemos
isso como $\pSat{v}{!A}$. Essa relação também pode ser definida por indução sobre
a estrutura de~$!A$, usando as funções de verdade dos conectivos
lógicos para definir, por exemplo, a satisfação de $!A \land !B$ em termos da
satisfação (ou não) de $!A$ e~$!B$.

Com base na relação de satisfação $\pSat{v}{!A}$ para sentenças,
podemos então definir as noções semânticas básicas de tautologia,
consequência e satisfatibilidade. !!^a{formula} é uma tautologia,
$\Entails !A$, se toda !!{valuation} a satisfaz, isto é,
$\pValue{v}(!A) = \True$ para toda~$\pAssign{v}$. Ela é consequência de um
conjunto de !!{formula}s, $\Gamma \Entails !A$, se toda !!{valuation} que
satisfaz todas as !!{formula}s de~$\Gamma$ também satisfaz~$!A$. E um
conjunto de !!{formula}s é satisfatível se alguma !!{valuation} satisfaz simultaneamente todas as
!!{formula}s que ele contém. Como as !!{formula}s são
definidas indutivamente, e a satisfação, por sua vez, é definida por indução
sobre a estrutura das !!{formula}s, podemos usar a indução para demonstrar
propriedades de nossa semântica e relacionar as noções semânticas
definidas.


\end{document}
